\section{Conclusion}
\label{sec:conclusion}

We have argued that Case-Based Reasoning in LLM workflow execution requires a reasoning
medium, and demonstrated that when such a medium exists, CBR is realized natively at
two levels without requiring additional machinery.

\textbf{Level~1} is immediate: every execution checkpoint is a suspended NormCode
runtime --- a self-contained, resumable case whose integrity is guaranteed by the scope
rule. The four CBR operations map directly onto Canvas interactions.

\textbf{Level~2} is the deeper contribution: each compiled NormCode plan is an abstract
case --- a runtime definition from which all executions of that task type are
instantiated. Four production-ready plans serve as an initial case library. The
compilation pipeline is itself a NormCode plan --- self-hosting compilations as
recursive case-based learning.

\textbf{Three measurable properties} follow structurally: C1, C2, C3. These properties
are absent from existing LLM workflow tools --- not because those tools lack features,
but because they lack the reasoning medium that would make these properties structural
rather than optional.

\textbf{For CBR.} NormCode contributes a new case type (the suspended runtime), a new
case quality guarantee (structural isolation as a formal precondition), and a concrete
deployed instance of two-level POCBR. The four production-ready plans and two case
studies provide initial evidence that the architecture operates as claimed.

\textbf{For LLM workflows.} The CBR framing clarifies what the checkpoint system is: a
case base from which past experience can be retrieved, reused, and improved. This
reframing suggests research investments --- learned similarity for automated fork
selection, competence-based case base maintenance, cross-plan reuse.

\textbf{Availability.} NormCode Canvas (v1.1.3) is available for download. The language,
compiler, and orchestrator are open. A live web demo of the PPT Agent is available. The
formal language specification is documented in \cite{anon2025}.

\textbf{Future work.}
\begin{enumerate}
  \item Learned similarity measures for automated fork target selection.
  \item Case base maintenance policies.
  \item Formal evaluation against baseline frameworks on established multi-step
    reasoning benchmarks.
  \item Cross-plan Level~2 retrieval --- adapting abstract cases across workflow domains.
\end{enumerate}
